\chapter{SSA Destruction and Machine Optimization}\label{ch:ssa-destruction-machine-opt}

Static single assignment simplifies the middle end but no machine executes it directly.
Hardware has no phi functions, no infinite temporaries, and no patience for redundant copies.
Some pass must translate SSA away, repair the seams with real moves, and then squeeze the machine code with low-level cleanups.
This chapter covers that exit ramp: correct SSA destruction followed by the machine optimizations that pay for the trouble.

\section{Leaving SSA Correctly}\label{sec:ch28-destruction}

SSA destruction replaces each phi function with moves that select the right incoming value along each path.
\term{SSA destruction} (also called deconstruction or out-of-SSA translation) eliminates phi functions while preserving the value flow they encoded.
\term{Phi congruence class} means the set of temporaries linked by phi functions that want the same location.
\term{Critical edge} means an edge from a block with several successors to a block with several predecessors, where a placed move would execute on the wrong paths.
\term{Parallel copy} means a set of moves that take effect simultaneously at a point, so that swapping two registers needs a temporary rather than sequential overwrites.
The naive rule of copying each phi input in its predecessor works until two phis in one block interact or an edge is critical.

Consider a join block with two phis that swap their inputs.
Sequential copies in the predecessors or the join block overwrite one value before reading the other.
Parallel-copy semantics fixes the order by reading all sources before writing any destination.
Critical edges force a different repair: the move cannot sit in the source because it would run for all successors, and it cannot sit in the destination because values from different predecessors would collide.
Splitting the edge with a fresh block gives the move a home that executes on exactly one path.
The snippet below shows the swap that defeats sequential copies.
The \code{phi} operations are \code{SSA constructs} with no direct machine equivalent.
\begin{irlisting}
; join block J with predecessors A and B
J: x2 = phi(x0 from A, x1 from B)
   y2 = phi(y0 from A, y1 from B)
; swap case: x0 is y1 and y0 is x1 in values
\end{irlisting}
After destruction with parallel copies, block \code{A} ends with \code{x2, y2 := x0, y0} taken atomically, and block \code{B} ends with \code{x2, y2 := x1, y1} atomically.
If allocation gives the sources and destinations crossed locations, a copy-sequencing pass expands the parallel copy through a temporary.
Two distinct simultaneously live results cannot share one register.
No use between the copy and the join observes an intermediate state.

Figure~\ref{fig:visual-copy-snapshot} separates the two phases of a parallel copy: read every old source, then write the destinations.
Listing~\ref{lst:visual-parallel-copy} implements that definition directly for eight simulated registers.
Its temporary array is a simple reference implementation; a machine-code sequencer can use fewer temporary locations by finding cycles, as discussed below.

\begin{figure}[htbp]
  \centering
  \begin{tikzpicture}[node distance=12mm and 16mm]
    \node[basic block,align=center] (r1) {old $r_1$\\1};
    \node[basic block,align=center,right=of r1] (r2) {old $r_2$\\2};
    \node[basic block,align=center,below=of r1] (s0) {saved[0]\\2};
    \node[basic block,align=center,below=of r2] (s1) {saved[1]\\1};
    \node[basic block,align=center,below=of s0] (n1) {new $r_1$\\2};
    \node[basic block,align=center,below=of s1] (n2) {new $r_2$\\1};
    \draw[cfg edge] (r2) -- (s0);
    \draw[cfg edge] (r1) -- (s1);
    \draw[dashed] ([xshift=-6mm,yshift=-5mm]s0.south west) --
      ([xshift=6mm,yshift=-5mm]s1.south east);
    \draw[cfg edge] (s0) -- (n1);
    \draw[cfg edge] (s1) -- (n2);
    \node[font=\sffamily\small,anchor=west,right=8mm of s1] {all reads finish before writes};
  \end{tikzpicture}
  \caption{Saving both sources makes a swap independent of destination-write order; the dashed boundary separates the read and write phases.}
  \label{fig:visual-copy-snapshot}
\end{figure}



\begin{listing}[htbp]
  \inputminted{c}{example-code/visual-labs/copy.c}
  \caption{Snapshotting source values preserves swaps and longer copy cycles without overwriting an unread source.}
  \label{lst:visual-parallel-copy}
\end{listing}

The header defines a \code{Copy} as a destination/source index pair and fixes \code{REGISTER\_COUNT} at eight.
Indices must lie in that range, and each destination may appear only once; repeated sources are permitted.
All validation and source reads finish before the first register write, so rejection leaves the register array unchanged.
For $m$ copies the two loops take $O(m)$ time; the fixed temporary arrays provide bounded storage.
The accompanying test checks every pair of byte-valued swap inputs, a three-register cycle, a repeated source, and rejection after an earlier valid copy.

Lost-copy and swap problems are the two classic failure modes.
The \term{lost-copy problem} drops a value when a phi result and its input share a location that an unrelated definition overwrites along one path.
The \term{swap problem} corrupts two phis whose sources and destinations cross, as above.
Both vanish when destruction uses parallel copies placed on non-critical edges plus systematic edge splitting.
Extensions for predicated code, hardware register pairs, and exception edges add constraints on where copies may sit, but the skeleton stays the same.
Figure~\ref{fig:ch28-critical} shows why the fresh block is needed.
Table~\ref{tab:ch28-destruction-fixes} maps each failure to its fix.

\begin{figure}[htbp]
\centering
\begin{tikzpicture}
\node[font=\sffamily\small] at (0.8,2.3) {Before};
\node[basic block] (a) at (0,1.5) {A};
\node[basic block] (c) at (2,1.5) {C};
\node[basic block] (j) at (1,-0.3) {J: phis};
\node[basic block] (z) at (-1,-0.3) {Z};
\draw[cfg edge] (a) -- node[right,font=\small] {critical} (j);
\draw[cfg edge] (a) -- (z);
\draw[cfg edge] (c) -- (j);
\begin{scope}[xshift=65mm]
\node[font=\sffamily\small] at (0.8,2.3) {After};
\node[basic block] (aa) at (0,1.5) {A};
\node[basic block] (cc) at (2,1.5) {C};
\node[basic block] (ss) at (0,0) {S: copies};
\node[basic block] (jj) at (1,-1.5) {J};
\node[basic block] (zz) at (-1.4,-1.5) {Z};
\draw[cfg edge] (aa) -- (ss);
\draw[cfg edge] (aa.west) -- ++(-6mm,0) -| (zz.north);
\draw[cfg edge] (ss) -- (jj);
\draw[cfg edge] (cc) -- (jj);
\end{scope}
\end{tikzpicture}
\caption{Splitting $A\to J$ puts its phi copies in $S$, so they run neither on $A\to Z$ nor on $C\to J$.}
\label{fig:ch28-critical}
\end{figure}

\begin{table}[htbp]
\caption{SSA-destruction hazards with their standard repairs.}
\label{tab:ch28-destruction-fixes}
\begin{tabular}{lll}
\toprule
Hazard & Symptom & Repair \\
\midrule
Critical edge & No single home for copy & Split the edge \\
Lost copy & Overwritten shared location & Parallel copy plus split \\
Swap & Crossed sources clobbered & Temporary in sequence \\
Dead phi & Useless copies survive & Dead-code pass after \\
\bottomrule
\end{tabular}
\end{table}

Destruction order interacts with coalescing and copy folding.
Merging a phi result with one of its inputs removes a copy but creates a congruence that later spilling must respect.
Over-eager merging reintroduces the lost-copy shape under register pressure.
Modern out-of-SSA schemes therefore keep the parallel copies explicit, let the allocator coalesce what stays colorable, and sequence whatever remains.
A final dead-code pass removes phis whose results are unused, which happens more often than newcomers expect after earlier simplification.
\begin{ccwarning}[Do not place copies by intuition]
Hand-placed moves on critical edges look right on the common path and corrupt the rare path.
Always split first, use parallel-copy semantics, and expand to sequential moves mechanically.
Test with swapped inputs explicitly.
\end{ccwarning}

\section{Out-of-SSA Translation}\label{sec:ch28-algorithm}

Algorithm~\ref{alg:ch28-out-of-ssa} gives the standard destruction pass used before or during register allocation.
It splits critical edges on demand, installs parallel copies, sequences them without breaking swaps, and cleans up.
The input is SSA with explicit phi functions and a correct control-flow graph.
The output is ordinary three-address or machine IR with no phis and equivalent value flow.
The algorithm preserves semantics because each copy executes exactly along the path its phi input described, and parallel semantics removes ordering hazards.

\begin{algorithm}[htbp]
\caption{Out-of-SSA translation with edge splitting}\label{alg:ch28-out-of-ssa}
\begin{pseudocode}
\Require SSA function with phi nodes, CFG $G$
\Ensure equivalent function with no phi nodes
\ForAll{critical edge $e$ that needs a phi copy}
\State split $e$ with a fresh block $S_e$
\EndFor
\ForAll{block $J$ with phi nodes}
\ForAll{predecessor $P$ of $J$}
\State build parallel copy $C_P$ mapping each phi input from $P$ to its result
\State place $C_P$ on edge $P\to J$: before $P$'s terminator if it has one successor, otherwise at $J$'s entry if it has one predecessor, or in the split block
\EndFor
\State delete the phi nodes in $J$
\EndFor
\State sequence each parallel copy with temporaries for swaps
\State run dead-code elimination on orphaned copies
\end{pseudocode}
\end{algorithm}

Sequencing a parallel copy needs its own small routine: repeatedly emit copies whose destination is not a remaining source, and break cycles with one fresh temporary per cycle.
With source-use counts and a queue of ready copies, sequencing is linear in the copy count; repeatedly scanning all remaining copies can be quadratic.
The whole pass costs $O(N + E + P)$, where $N$ is the instruction count, $E$ is the edge count, and $P$ is the total phi-input count, plus the cost of the edge splits it creates.
Memory behavior is predictable because each phi input becomes exactly one parallel-copy element and then at most two sequential moves.
The pass runs once per function in the common pipeline, before allocation, with a small post-allocation repair for copies the allocator could not coalesce.
% TODO(cite): need a primary reference for parallel-copy SSA destruction.

\section{Worked Example}\label{sec:ch28-worked}

We destroy two phis by hand on a diamond with a critical edge.
Blocks \code{A} and \code{C} feed join block \code{J}, and \code{A} also branches to an unrelated block \code{Z}, so edge \code{A -> J} is critical.
\begin{irlisting}
A: x0 = const 1; y0 = const 2; br J or Z
C: x1 = const 3; y1 = const 4; jmp J
J: x2 = phi(x0 from A, x1 from C)
   y2 = phi(y0 from A, y1 from C)
   use x2, y2
\end{irlisting}
The second phi takes \code{y0} from \code{A} and \code{y1} from \code{C}.
The values alone do not create a swap; a swap arises when register assignment crosses the source and destination locations.
Step one splits edge \code{A -> J} with fresh block \code{S}.
Step two installs parallel copies: at the end of \code{S} put \code{x2, y2 := x0, y0}, and at the end of \code{C} put \code{x2, y2 := x1, y1}.
Step three deletes both phis from \code{J}.
Step four sequences the copies.
The \code{C} copy has disjoint sources and destinations in this numbering, so it becomes two plain moves in either order.
Suppose \code{x0} occupies register $r_1$, \code{y0} occupies $r_2$, and the destinations \code{x2} and \code{y2} occupy $r_2$ and $r_1$, respectively.
The required parallel copy is $r_2\gets r_1, r_1\gets r_2$.
Save $r_1$ in a temporary, copy $r_2$ into $r_1$, and copy the temporary into $r_2$.
Starting with $(r_1,r_2)=(1,2)$ leaves $(r_1,r_2)=(2,1)$, so \code{x2} is $1$ and \code{y2} is $2$ as required.

Verify path by path.
Along \code{A -> S -> J}, the join observes \code{1} and \code{2} in the right slots with no leakage into the \code{A -> Z} path, because the copy sits in \code{S} rather than in \code{A}.
Along \code{C -> J}, the join observes \code{3} and \code{4}.
Delete-check: if \code{y2} were unused after \code{J}, dead-code elimination would remove its copy elements and leave only the \code{x2} moves.
Count the final moves: two on the \code{C} path, three (including the temporary) on the swap path.
That count is minimal for the chosen assignment; a different register assignment could drop the temporary, which is exactly the coalescing decision left to the allocator.

\section{Machine Optimizations After SSA}\label{sec:ch28-machine-opt}

Once phis are gone, a family of machine-level cleanups reclaims the copies and sharpens the code.
Copy propagation and \term{peephole optimization} remove redundant moves, fold flag-setting comparisons into branches, and replace generic sequences with target idioms.
\term{Branch tuning} shortens jumps, aligns hot targets, removes unreachable blocks, and converts predictable branches into conditional moves or predication where profitable.
\term{Machine-level dead-code elimination} drops flag computations, unused condition codes, and copies whose results never reach a use.
\term{Stack-slot compaction} and frame-index elimination shrink the frame that allocation pessimized.
Each of these is a small pattern rewrite, but together they decide whether the backend output looks hand-tuned or machine-generated.

The ordering rule is simple: destroy SSA, allocate registers, then optimize the machine code with full knowledge of real registers and frame offsets.
Running peephole logic before allocation wastes effort on copies the allocator will rewrite anyway.
Running branch tuning after emission is too late to fix layout.
A compact post-allocation loop of copy cleanup, dead-code removal, branch shortening, and block placement captures most of the gain in linear or near-linear time per pass.
Profile-guided block placement deserves a mention because it moves cold code out of line, shortens hot fall-through paths, and improves instruction-cache density with no semantic risk beyond edge-count fidelity.

Test destruction and machine cleanup together rather than separately.
Random CFG fuzzers with phi swaps and critical edges catch lost-copy regressions that unit tests miss.
Check exception edges and debug ranges after edge splitting, since fresh blocks shift line mappings.
Differential-test optimized machine code against unoptimized machine code on input corpora, including NaN, signed zero, and overflow corners for the folded flag sequences.
Background on backend pass structure and low-level cleanup appears in~\cite{morgan-1998-optimizing} and see~\cite{appel-1998-modern}.
Validation discipline for these late semantics-preserving rewrites follows see~\cite{safonov-2010-trustworthy}.
Language facts about observable behavior that late rewrites must preserve are stated in~\cite{wg14-2011-c11}.

\section{Exercises}\label{sec:ch28-exercises}

\begin{enumerate}
\item A join block has one phi fed by two predecessors, and one incoming edge is critical. \label{ex:ch28-split-place} Split the edge, place the copies, and show why placing the copy in the source block misexecutes.
\item Two phis in one join swap their inputs along one path. \label{ex:ch28-swap-seq} Write the parallel copy, sequence it with a temporary, and count the final moves.
\item A post-allocation peephole sees \code{cmp r, 0} followed by \code{br\_eq L}. \label{ex:ch28-peephole} Give the folded target idiom, state its preconditions, and explain what the dead-code pass removes if the flags are otherwise unused.
\end{enumerate}

\section{Further Reading}\label{sec:ch28-reading}

Low-level pass structure, peephole design, and branch handling are covered in~\cite{morgan-1998-optimizing}.
SSA handling and backend organization from a complementary view appear in~\cite{appel-1998-modern}.
Trust arguments for late semantics-preserving transformations appear in~\cite{safonov-2010-trustworthy}.
Observable-behavior rules that constrain late rewrites are defined in~\cite{wg14-2011-c11}.
